% !TeX program = xelatex
% UTF-8 transcription; source wording, abbreviations, and doses are retained.
% images/pic01.jpg = 20261012_190655602.jpg
% images/pic02.jpg = 20261012_190702951.jpg
% images/pic03.jpg = 20261012_190709672.jpg
% images/pic04.jpg = 20261012_190716735.jpg
% images/pic05.jpg = 20261012_190724572.jpg
% images/pic06.jpg = 20261012_190731623.jpg
% images/pic07.jpg = 20261012_190737611.jpg
% images/pic08.jpg = 20261012_190742560.jpg
% Standalone XeLaTeX document, also usable via input in a XePersian document.
\ifdefined\settextfont
  \def\HFrEFStandalone{0}
\else
  \def\HFrEFStandalone{1}
  \documentclass[12pt,a4paper]{article}
  \usepackage[margin=2cm]{geometry}
  \usepackage{amsmath,amssymb}
  \usepackage{array}
  \usepackage{xepersian}
  \IfFontExistsTF{Amiri}{\settextfont{Amiri}}{\settextfont{FreeSerif}}
  \setlatintextfont{FreeSerif}
\fi
\providecommand{\en}[1]{\lr{#1}}
\newcommand{\HFrEFpage}[2]{%
  \clearpage
  \section*{موضوع: براونوالد ـ ۵۰ ـ منج \en{#2}}
  {\footnotesize تاریخ: \en{/ /}\quad صفحهٔ #1\par}}
\ifnum\HFrEFStandalone=1 \begin{document}\fi

\HFrEFpage{۱}{HFrEF}
\begin{itemize}
\item اپروچ به بیمار \en{HFrEF}: یک طیف است و ۴ استیج مختلف دارد.
\item \en{Stage A}: بیمار در ریسک بالای ایجاد \en{HF} است ولی در حال حاضر اختلال استراکچرال قلبی ندارد و علائمی از \en{HF} نیز ندارد.
  \begin{itemize}
  \item بیمار \en{DM} و \en{HTN}.
  \end{itemize}
\item \en{Stage B}: بیمار اختلال استراکچرال قلبی دارد ولی علامت بالینی ندارد.
  \begin{itemize}
  \item بیماری که \en{MI} قبلی داشته و الان افت \en{EF} دارد.
  \end{itemize}
\item \en{Stage C}: اختلال استراکچرال دارد و علامت بالینی نیز دارد.
\item \en{Stage D}: علی‌رغم درمان همچنان علامت‌دار هستند و در \underline{لیست پیوند قلب}.
\end{itemize}
\subsection*{مدیریت \en{Stage A}}
\begin{itemize}
\item جلوگیری از ایجاد \en{HF}.
\item کنترل \en{DM}، \en{HTN}، هایپرلیپیدمی.
\item با اندیکاسیون‌ها نیز \en{ACEI} بدهیم (\en{ARB} هم می‌شود).
  \begin{itemize}
  \item بیماری آترواسکلروتیک عروقی، \en{DM}، \en{HTN} + ریسک فاکتورهای قلبی عروقی.
  \end{itemize}
\item در صورت نیاز شروع استاتین.
\item بیماران با سابقهٔ خانوادگی کاردیومیوپاتی و بیماران با دریافت داروهای کاردیوتوکسیک باید غربالگری دوره‌ای شوند از نظر \en{HF}: اکو.
\end{itemize}
منج بیمارانی که \en{HF} دارند \en{$\pm$} علامت \en{$\Longleftarrow$}

\HFrEFpage{۲}{HFrEF}
\begin{itemize}
\item سندرم تاکاتسوبو: اختلال عملکرد برگشت‌پذیر \en{LV} به دنبال استرس روحی که با \en{CP}، ادم ریه و شوک کاردیوژنیک همراه است و بیماری کرونری هم ندارد.
\item به دنبال یک اختلال عملکرد \en{LV} در زمینهٔ آسیب میوکارد، بیمار ممکن است تا چندین ماه بدون علامت باشد.
\item برای بیمارانی که اختلال عملکرد بطن چپ دارند ولی بدون علامت هستند (\en{Class I}) هدف درمان کاهش سرعت پیشرفت بیماری با بلاک کردن سیستم نوروهمورال و جلوگیری از ریمودلینگ قلب است.
\item برای بیمارانی که علامت‌دار شده‌اند (\en{Class II to IV}) هدف درمان:
  \begin{itemize}
  \item رفع احتباس مایعات.
  \item کاهش دیس‌ابیلیتی بیماران.
  \item کاهش خطر پیشرفت بیماری و مرگ.
  \end{itemize}
\item بیماری‌های همزمان که حتماً در درمان \en{HF} باید کنترل شوند:
  \begin{itemize}
  \item \en{HTN}، آنمی، \en{Sleep-disordered breathing}.
  \end{itemize}
\end{itemize}

\HFrEFpage{۳}{HFrEF}
\begin{itemize}
\item مصرف سیگار باید قطع شود.
\item مصرف الکل: ۲ درینک در روز برای آقایان و ۱ درینک برای خانم‌ها.
\end{itemize}
\subsection*{ریسک فاکتورهای دکامپنسیشن \en{HF}}
\begin{itemize}
\item مصرف نامنظم داروها، \en{MI}، آریتمی، عفونت، آنمی، بارداری و …
\item شروع داروهایی که می‌توانند علائم \en{HF} را تشدید کنند:
  \begin{itemize}
  \item \en{BB}، \en{NSAID}، تیازولیدیندیون‌ها، آنتی‌بادی ضد \en{TNF}.
  \item داروهای آنتی‌آریتمی (تمام) کلاس \en{I}، سوتالول (کلاس \en{III}).
  \end{itemize}
\end{itemize}
\subsection*{دوز دیورتیک‌ها در بیمار \en{HF}}
\begin{itemize}
\item دیورتیک لوپ.
\end{itemize}
\begin{center}
\begin{tabular}{rccc}
دارو & دوز شروع & حداکثر دوز & طول اثر \\
فوروزماید & \en{20--40 mg daily/BD} & \en{600 mg} & \en{6--8 h} \\
\end{tabular}
\end{center}
\begin{itemize}
\item دیورتیک تیازید.
\end{itemize}
\begin{center}
\begin{tabular}{rccc}
دارو & دوز شروع & حداکثر دوز & طول اثر \\
\en{HCTZ} & \en{25 mg daily/BD} & \en{200 mg} & \en{6--12 h} \\
متولازون & \en{2.5--5 mg daily} & \en{5 mg} & \en{12--24 h} \\
\end{tabular}
\end{center}
\begin{itemize}
\item پتاسیم اسپیرینگ.
\end{itemize}
\begin{center}
\begin{tabular}{rccc}
دارو & دوز شروع & حداکثر دوز & طول اثر \\
آمیلوراید & \en{5 mg daily} & \en{20 mg} & \en{24 h} \\
\end{tabular}
\end{center}

\HFrEFpage{۴}{HF}
\subsection*{فعالیت فیزیکی در بیمار \en{HF}}
\begin{itemize}
\item فعالیت فیزیکی شدید توصیه نمی‌شود.
\item فعالیت روتین و ورزش سبک به بیماران کلاس \en{I} تا \en{III} توصیه می‌شود.
\item برای بیمار، دویدن سبک، پیاده‌روی و دوچرخهٔ ثابت منعی ندارد.
\item تا ۶ هفته بعد از یک حادثهٔ قلبی عروقی، با ملاحظه، نباید ورزش کند.
\item سایر کنتراندیکاسیون‌های ورزش:
  \begin{itemize}
  \item وجود یک دیوایس قلبی که اجازه ندهد \en{HR} بالا برود.
  \item وجود یک آریتمی \en{Significant} یا ایسکمی در تست ورزش.
  \end{itemize}
\end{itemize}
\subsection*{رژیم غذایی در بیمار \en{HF}}
\begin{itemize}
\item در همهٔ بیماران \en{HF} سدیم روزانه زیر \en{3 gr} باشد؛ در \en{HF} متوسط تا شدید زیر \en{2 gr}.
\item محدودیت مایع لزومی ندارد مگر بیمار هایپوناترمیک باشد (\en{$Na<130$}).
  \begin{itemize}
  \item زیر \en{2 lit} روزانه.
  \end{itemize}
\item به دلیل کاشکسی قلبی در بیماران \en{HF} پیشرفته، باید رژیم غذایی پرکالری باشد.
\item مصرف مکمل‌های تغذیه‌ای شواهدی ندارد و کمک‌کننده نیست.
\end{itemize}
\subsection*{منج احتباس مایعات در \en{HF}}
\begin{itemize}
\item دیژیتال‌ها و \en{ACEI} با دوز پایین ترشح \en{Na} در ادرار را زیاد می‌کنند.
\item ولی برای بیماران علامت‌دار حتماً دیورتیک نیز لازم است.
\end{itemize}

\HFrEFpage{۵}{HF}
\begin{itemize}
\item لوپ دیورتیک ترشح \en{Na} را ۲۰–۲۵\% بالا می‌برد.
\item تیازید ترشح \en{Na} را ۵–۱۰\% بالا می‌برد.
\item تیازید اثر خود را در \en{$GFR<40$} از دست می‌دهد.
\item دیورتیک لوپ با \en{Cl} بر سر اتصال به \en{NKCC2} در غشای \en{Apical} سلول‌ها در بخش \en{thick Ascending} لوپ هنله رقابت می‌کند.
\item فراهمی زیستی آن ۴۰–۷۰\% در خوراکی است (فوروزماید).
\item تیازیدها می‌توانند هایپوناترمی ایجاد کنند.
  \begin{itemize}
  \item می‌توانند هایپرکلسمی خفیف ایجاد کنند.
  \end{itemize}
\item در طول زمان هایپومنیزیمی ایجاد می‌کنند.
\item هایپراوریسمی می‌توانند ایجاد کنند.
\item یادداشت حاشیه: \en{$K\downarrow$}، \en{$Na\downarrow$}، \en{$Mg\downarrow$}، \en{$Ca\uparrow$} (دایره دور \en{Ca}).
\item مینرالوکورتیکوئیدها / آنتاگونیست‌های گیرندهٔ مینرالوکورتیکوئید:
  \begin{itemize}
  \item \en{MRAs}: آلدوسترون؛ حفظ آب و نمک و دفع \en{$K^+$} و \en{$H^+$}.
  \item این اتفاق را مهار می‌کنند.
  \item عارضهٔ مهم آن‌ها هایپرکالمی.
  \end{itemize}
\item دیورتیک‌های پتاسیم اسپیرینگ:
  \begin{itemize}
  \item \en{Amiloride}، \en{Triamterene}.
  \item افزایش خفیف ترشح \en{NaCl}.
  \end{itemize}
\item مهارکننده‌های کربونیک انهیدراز: استازولاماید.
  \begin{itemize}
  \item مهار بازجذب \en{$NaHCO_3$} در توبول پروگزیمال.
  \item در درمان آلکالوز متابولیک در زمینهٔ درمان با سایر دیورتیک‌ها.
  \item عارضه: اسیدوز متابولیک + هایپوکالمی شدید.
  \end{itemize}
\end{itemize}

\HFrEFpage{۶}{HF}
\begin{itemize}
\item \en{SGLT2-I}: امپاگلیفلوزین.
  \begin{itemize}
  \item گیرندهٔ \en{SGLT2} مسئول بازجذب ۹۰\% گلوکز و همچنین بازجذب سدیم در توبول پروگزیمال است.
  \item سطح اوریک‌اسید پلاسما را نیز در بیماران ۱۰–۱۵\% کاهش می‌دهد.
  \end{itemize}
\item در بیمار \en{HF} دیورتیک همیشه باید همراه با آنتاگونیست نوروهمورال باشد.
\item دوز شروع فوروزماید در بیمار با \en{GFR} نرمال \en{40 mg} است.
\item وقتی بیمار در زمینهٔ درمان با دیورتیک خشک شد، باید وزن او را یادداشت کنیم.
\item عارضهٔ درمان با فوروزماید:
  \begin{itemize}
  \item اختلال الکترولیتی، کاهش حجم مایعات، تشدید آزوتمی.
  \end{itemize}
\item سطح پتاسیم بیمار \en{HF} باید در رنج ۴–۵ حفظ شود.
\end{itemize}
\subsection*{جلوگیری از پیشرفت \en{HF} با دارو}
\begin{itemize}
\item همهٔ بیماران \en{HF} باید \en{ACEI} دریافت کنند.
  \begin{itemize}
  \item یادداشت بالای \en{HF}: \en{HFrEF ($LVEF<40\%$)}.
  \item جلوگیری از ریمودلینگ قلبی.
  \end{itemize}
\end{itemize}

\HFrEFpage{۷}{HF}
\begin{itemize}
\item قبل از شروع \en{ACEI} باید دوز فوروزماید را به اپتیمم برسانیم.
\item با دوز کم شروع می‌کنیم و هر ۳–۵ روز دوز دارو را ۲ برابر کنیم.
\item در بیماران \en{Stable} قبل از آنکه دوز \en{ACEI} را حداکثر کنیم می‌توانیم \en{BB} را نیز شروع کنیم.
\item ۱–۲ هفته پس از شروع \en{ACEI}، عملکرد کلیه، پتاسیم و \en{BP} را چک.
\end{itemize}
\subsection*{داروها جلوگیری از پیشرفت \en{HF} دهند}
\begin{itemize}
\item \en{ACEI}
  \begin{itemize}
  \item کاپتوپریل: \en{$6.25\,mg\ TDS\ \longrightarrow\ 50\,mg\ TDS$}.
  \item انالاپریل: \en{$2.5\,mg\ BD\ \longrightarrow\ 10\,mg\ BD$}.
  \end{itemize}
\item \en{ARB}
  \begin{itemize}
  \item والسارتان: \en{$40\,mg\ BD\ \longrightarrow\ 160\,mg\ BD$}.
  \item لوزارتان: \en{12.5--25 daily} \en{$\longrightarrow$} \en{50 mg daily}.
  \end{itemize}
\item \en{ARNI}
  \begin{itemize}
  \item ساکوبیتریل / والسارتان: \en{24/26 BID} \en{$\longrightarrow$} \en{97/103 BID}.
  \end{itemize}
\item \en{MRA}
  \begin{itemize}
  \item اسپیرونولاکتون: \en{12.5--25 daily} \en{$\longrightarrow$} \en{25--50 daily}.
  \item اپلرنون: \en{25 mg daily} \en{$\longrightarrow$} \en{50 mg daily}.
  \end{itemize}
\end{itemize}

\HFrEFpage{۸}{HF}
\begin{itemize}
\item به بیماران \en{ARB} می‌دهیم که دچار عارضهٔ \en{ACEI} شده‌اند.
  \begin{itemize}
  \item سرفه، راش پوستی، آنژیوادم.
  \item \en{ARB}: \en{COR II B}.
  \end{itemize}
\item در بیماران \en{HFrEF} مزمن کلاس \en{II} و \en{III} با \en{COR I} توصیه می‌شود که \en{ARNI} جایگزین \en{ACEI} شود.
\end{itemize}
\subsection*{بتابلاکر در \en{HF}}
\begin{itemize}
\item وقتی بیمار یوولمیک است شروع می‌کنیم و تدریجی دوز را بالا می‌بریم.
\item \en{COR: I} برای \en{HFrEF}، \en{$EF<40$}، چه علامت دارد چه بدون علامت.
\item تنها این سه \en{BB} کاربرد دارند:
  \begin{itemize}
  \item بیزوپرولول (\en{$\beta_1$}).
  \item متوپرولول سوکسینات (\en{$\beta_1$}).
  \item کارودیلول (\en{$\beta_2,\beta_1+\alpha_1$}).
  \end{itemize}
\item با فاصلهٔ زمانی هر ۲ هفته باید دوز را بالا ببریم.
\item عوارض شروع درمان ظرف ۳–۵ روز با افزایش وزن نمایان می‌شود.
  \begin{itemize}
  \item دوز دیورتیک را بالا می‌بریم.
  \end{itemize}
\end{itemize}

% Notebook subject/date labels and handwritten page numbers are retained.
% The printed notebook emblem contains no text.
\ifnum\HFrEFStandalone=1
  \let\HFrEFEnd\enddocument
\else
  \let\HFrEFEnd\relax
\fi
\HFrEFEnd
